% Plain article layout. The generated overleaf/main.tex embeds this preamble.
\documentclass[11pt,letterpaper]{article}
\usepackage[T1]{fontenc}
\usepackage{lmodern}
\usepackage[margin=1in]{geometry}
\usepackage{microtype}
\usepackage{amsmath,amssymb}
\usepackage{booktabs}
\usepackage{caption}
\usepackage{placeins}
\usepackage{fancyhdr}
\usepackage[hidelinks]{hyperref}

\setlength{\parindent}{1.5em}
\setlength{\parskip}{0pt}
\setlength{\headheight}{14pt}
\setlength{\emergencystretch}{2em}
\renewcommand{\arraystretch}{1.12}
\captionsetup{font=small,labelfont=bf,skip=7pt}
\pagestyle{fancy}
\fancyhf{}
\fancyhead[L]{\small Ovarian cancer data preparation}
\fancyfoot[C]{\thepage}
\renewcommand{\headrulewidth}{0.4pt}
\fancypagestyle{plain}{\fancyhf{}\fancyfoot[C]{\thepage}\renewcommand{\headrulewidth}{0pt}}
\providecommand{\tightlist}{\setlength{\itemsep}{0pt}\setlength{\parskip}{0pt}}
\hypersetup{pdftitle={Preparing an Ovarian Cancer Dataset},pdfsubject={BIOSTAT 721 Project 1}}

\title{Preparing an Ovarian Cancer Dataset}
\author{BIOSTAT 721 Project 1}
\date{}

\begin{document}
\maketitle
\vspace{-1.5em}
\section{Introduction}\label{introduction}

Women with active ovarian cancer attend chemotherapy appointments every
two to three weeks. Visits to the emergency department or unanticipated
hospital admissions between appointments are a concern for their
physicians. Predicting these events at the end of a chemotherapy
appointment first requires patient and encounter records that can be
used together.

According to the requirement, this project prepares such a dataset from
the two files supplied for BIOSTAT 721 Project 1. It combines patient
characteristics with encounter measurements, addresses missing and
implausible values, and summarizes white blood cell count (WBC) and
patient characteristics.

\section{Data}\label{data}

The patient file contains 50 rows and 8 variables, with one row per
patient. Its fields record the medical record number (MRN), date of
birth, race, ethnicity, financial class, hypertension, congestive heart
failure, and diabetes. The encounter file contains 550 rows and 7
variables, with one row per encounter. It records MRN, encounter date
and type, temperature, distress score, WBC, and body mass index (BMI).
The encounters span January 02, 2014 to September 01, 2019 and comprise
497 office visits and 53 hospital encounters.

Patient records were joined to encounter records by MRN. Each patient
identifier is unique in the patient file, and every encounter matches a
patient. Both files have no exact duplicate rows, and every combination
of MRN and encounter date is unique. The base analytic dataset therefore
contains 550 rows and 14 variables. Each row represents one encounter
for one of the 50 patients, with that patient's characteristics
attached. Adding the WBC category produces a final dataset with 550 rows
and 15 variables.

\section{Cleaning}\label{cleaning}

Dates were analyzed as calendar dates, and measurement fields as numeric
values. WBC is measured in \(10^3\) cells per \(\mu\)L. According to the
assignment requirement, the report treats any value below 1 as below the
detection limit and make it to be replaced by 1. This rule affected 14
encounters. It compresses the lower tail of the continuous WBC
distribution. Each affected observation remains below 3.2 after
replacement, so its WBC category is unaffected.

The 6 BMI entries recorded as \(-989\) were set to missing because
non-positive values cannot represent BMI. The remaining observed BMI
values range from 12.16 to 45.91. A birth date of July 22, 1889, implies
an age of 124 to 129 years across one patient's 17 encounters. An age
above 120 at an observed encounter was used as a project screening rule,
and this patient's date of birth was set to missing pending source
verification. Temperature values range from 96.26 to 103.20, but their
unit is not recorded. These temperatures and the positive BMI values
were retained but further check may still needed.

BMI was missing in 86 encounters (15.6\%) before cleaning and 92
(16.7\%) afterward. Distress score was missing in 53 encounters (9.6\%).
All 53 hospital encounters lack both BMI and distress score after
cleaning. Among the 497 office visits, 39 (7.8\%) lack BMI and all have
a distress score. WBC has no missing values. The suspect birth date
creates missing DOB for one patient and 17 encounter rows. All 550
encounters were retained, with missing values represented explicitly.
Original values and the reasons for each change are recorded in the
cleaning log.

\section{Results}\label{results}

Table \ref{tab:patients} describes the 50 unique patients. It was
calculated from the patient dataset so that each person contributes once
to each characteristic. All six characteristics are complete, and every
percentage uses 50 patients as its denominator. Most patients are White
(74.0\%) and non-Hispanic (94.0\%). Medicare is recorded for 29 patients
(58.0\%). Hypertension is recorded for 20 patients, congestive heart
failure for 5, and diabetes for 2.

\begin{table}[!htbp]
\centering
\caption{Patient characteristics ($N = 50$).}\label{tab:patients}

\begin{tabular}{llr}
\toprule
Characteristic & Category & n (\%)\\
\midrule
Race & White & 37 (74.0)\\
 & Black & 10 (20.0)\\
 & Other & 3 (6.0)\\
Ethnicity & non-Hispanic & 47 (94.0)\\
 & Hispanic & 3 (6.0)\\
Financial class & Medicare & 29 (58.0)\\
 & Private & 21 (42.0)\\
Hypertension & No & 30 (60.0)\\
 & Yes & 20 (40.0)\\
Congestive heart failure & No & 45 (90.0)\\
 & Yes & 5 (10.0)\\
Diabetes & No & 48 (96.0)\\
 & Yes & 2 (4.0)\\
\bottomrule
\end{tabular}

\par\vspace{5pt}\begin{minipage}{0.90\textwidth}\small
Percentages use all 50 patients. Each characteristic has no missing values.
\end{minipage}
\end{table}

After applying the detection-limit rule, WBC was classified as Low for
values below 3.2, Normal for values from 3.2 to 9.8 inclusive, High for
values above 9.8, and Not Taken for missing measurements. Table
\ref{tab:wbc} presents the counts in that order, using all 550
encounters as the denominator. WBC is Normal in 436 encounters (79.3\%),
Low in 54 (9.8\%), and High in 60 (10.9\%). The Not Taken category is
included with a count of zero.

\begin{table}[!htbp]
\centering
\caption{WBC categories across encounters ($N = 550$).}\label{tab:wbc}

\begin{tabular}{llr}
\toprule
Category & Definition & n (\%)\\
\midrule
Low & Below 3.2 & 54 (9.8)\\
Normal & 3.2 to 9.8, inclusive & 436 (79.3)\\
High & Above 9.8 & 60 (10.9)\\
Not Taken & Missing & 0 (0.0)\\
Total & All encounters & 550 (100.0)\\
\bottomrule
\end{tabular}

\par\vspace{5pt}\begin{minipage}{0.90\textwidth}\small
WBC is measured in $10^3$ cells per $\mu$L. Percentages use all 550 encounters.
\end{minipage}
\end{table}

\FloatBarrier

\section{Discussion}\label{discussion}

The final dataset brings patient characteristics and encounter
measurements together with explicit cleaning rules and traceable
changes. The patient and encounter summaries describe different units of
observation. Repeated encounters from the same person may be related and
need further check for potential connection among datas. The
concentration of missing BMI and distress scores in hospital encounters
also requires attention.

Additional information is needed to define the proposed prediction task
as well. Specifically, The label Hospital Encounter does not separately
identify emergency department visits, planned admissions, and
unanticipated admissions. A valid outcome window would require the
timing and nature of hospital events and the next chemotherapy
appointment. Measurement availability times are also absent.

The analysis used R 4.6.0, with base R for data preparation. This report
was generated in R Markdown using rmarkdown 2.32 and knitr 1.51, and
converted to LaTeX with Pandoc 3.8.3. The assignment instructions and
the two supplied CSV files are the sources for the analysis.

\section{AI Use Disclosure}\label{ai-use-disclosure}

Microsoft Copilot equipped with ChatGPT 5.6 was used for checking if the
task was completed and strictly follow the project requirement. It
pointed out the mistakes that I did not follow the requirement of ``(For
WBC data) Any count \textless1 in our data is under the limit of
detection. Convert any value \textless1 to 1'' and did not use RMarkdown
or Quarto to finish the report, instead runs a simple R script and write
the report with latex manually with the analyzed result. I manually fix
the first part of Rscript of the WBC conversion, and AI assisted in
teaching me to write a R markdown report template, using function to
automatically put analyzed result into the template, and compile it via
latex. All AI result is checked manually.
\end{document}
